\section{Implementation}\label{sec:impl}

We implemented $\Pieval$ of \cref{sec:pcs}, compiled as in \cref{sec:proofs:bcs} with salted Merkle trees and Fiat--Shamir, as a self-contained Rust library, \texttt{kbfold} (about $2{,}200$ lines including tests and benchmarks; its dependencies are the BLAKE3 hash function and the operating system's random number generator, through the \texttt{getrandom} crate).
The code follows the paper line by line; this section records the concrete choices, and \cref{sec:exp} reports measurements.
The library also has a coefficient-form mode, used only as the baseline of \cref{sec:exp}.

\paragraph{Index conventions.}
The paper numbers bits, variables and rounds from $1$; the code, like most Rust code, numbers them from $0$.
Bit $b_k$ of the paper is bit $k-1$ of the code (\texttt{(b >> (k-1)) \& 1}), the variable $x_k$ and the coordinate $z_k$ are entries $k-1$ of the corresponding arrays, and the round polynomial $s_j$ and the challenge $r_j$ are stored at index $j-1$, while the oracle $w_j$ (sent in round $j+1$) is stored at index $j$.
Tables $f : \iv{0}{N-1} \to \Fq$ are arrays of length $N$ indexed identically in both.

\subsection{Fields}\label{sec:impl:fields}

\paragraph{Base field.}
Tables and the initial codeword $w_0$ live over the Goldilocks field $\mathbb{F}_p$, $p = 2^{64} - 2^{32} + 1$.
The largest power of $2$ dividing $p - 1$ is $2^{32}$, so smooth domains of order up to $2^{32}$ exist (\cref{lem:power}).
We take $L = \langle \omega \rangle$ with $\omega = 7^{(p-1)/M}$, where $7$ generates $\mathbb{F}_p^{\times}$: since $p - 1 = 2^{32}\cdot 3 \cdot 5 \cdot 17 \cdot 257 \cdot 65537$, this amounts to $7^{(p-1)/\ell'} \ne 1$ for each of the six prime factors $\ell'$, which the test suite checks.
Multiplication uses the standard reduction of a $128$-bit product modulo $p$ based on $2^{64} \equiv 2^{32}-1$ and $2^{96} \equiv -1$; elements are kept in canonical form.

\paragraph{Extension fields.}
The code is generic in the challenge field $\F$ of \cref{sec:pcs:params}, with two instances: the quadratic extension $\K = \mathbb{F}_p[\iota]/(\iota^2 - 7)$, of size $p^2 \approx 2^{128}$, for the classical parameters of \cref{rem:params} and \cref{sec:exp}; and the quartic extension $\mathbb{F}_p[\iota]/(\iota^4 - 7)$, of size $p^4 \approx 2^{256}$, for the post-quantum parameters of \cref{rem:pq-params}.
The challenges $r_j$, the evaluation point $z$, the value $v$, the sumcheck tables and all folded words $w_j$ ($j \ge 1$) live in $\F$, while tables and $w_0$ have entries in $\mathbb{F}_p$.
The polynomial $\iota^2 - 7$ is irreducible because $7$, a generator of $\mathbb{F}_p^\times$, is a non-square; $\iota^4 - 7$ is irreducible by \cref{rem:pq-params}.
The quartic extension is implemented as $\K[\upsilon]/(\upsilon^2 - \iota)$, where $\iota$ is the generator of $\K$; then $\upsilon^4 = 7$, and $\upsilon$ is the element written $\iota$ in $\mathbb{F}_p[\iota]/(\iota^4 - 7)$. Elements of $\F$ are written in the basis $1, \iota, \dots, \iota^{e-1}$, $e \in \{2, 4\}$ (the bijection $\mathrm{num}$ of \cref{sec:sound:pq}).
Multiplication in $\K$ uses three base-field multiplications (Karatsuba) and one multiplication by the constant $7$.
Mixing base-field words with extension-field challenges is only needed in the first fold, which the code handles generically.

\subsection{Encoding and folding}\label{sec:impl:enc}

\paragraph{Commit.}
$\Enc(f)$ (\cref{def:encoding}) is computed in three steps:
the kernel-to-monomial transform $\mathrm{coef}(U_f) = C^{\otimes n} f$ of \cref{cor:transforms}, applied in place one tensor factor at a time as $(x, x') \mapsto (x', x + x')$ (exactly $\tfrac{n}{2}N$ additions, no multiplications);
zero-padding to length $M$;
and an iterative radix-$2$ NTT producing $w_0(\omega^i)$ in natural order.

\paragraph{Domain layout.}
With natural ordering, $L_j = \langle \omega^{2^j} \rangle$ is stored as $(\omega^{2^j i})_{i \in \iv{0}{M_j-1}}$.
The fibre of index $i \in \iv{0}{M_{j+1}-1}$ is the pair of positions $(i, i + M_{j+1})$, since $\omega^{2^j M_{j+1}} = -1$, and its image under squaring is position $i$ of $L_{j+1}$.
Folding therefore reads the two halves of the array and writes one array of half size, with no permutation.

\paragraph{Fold.}
Both prover and verifier evaluate \eqref{eq:fold-explicit} in the line form of \cref{lem:line}.
With $a = w(\xi)$ and $a' = w(-\xi)$, they compute $a_e = w_e(\xi^2) = (a+a')/2$ and $a_o = w_o(\xi^2) = (a-a')/(2\xi)$, and then $\fold_r(w)(\xi^2) = (a_o - a_e) + r\,(2a_e - a_o)$ (\cref{lem:line}).
This costs one multiplication by $r$ in $\K$ and two multiplications by base-field constants.
The inverses $(2\xi)^{-1} = 2^{-1}\omega^{-2^j i}$ are read from a single table of $\omega^{-i}$, $i < M/2$, shared by all levels.
The same function is used by the prover on whole words and by the verifier on single fibres, which rules out the parametrisation mismatch of \cref{rem:pitfall} by construction.

\paragraph{Sumcheck.}
The prover keeps the tables $f_{j-1}$ and $\chi_{j-1}$ of \cref{sec:pcs:protocol} as arrays over $\K$.
It sends $s_j(0)$, $s_j(1)$ and $s_j(2)$, computed in one pass over the pairs of entries at positions $2c$ and $2c+1$ of both tables, using
\[
  s_j(2) \;=\; \sum_{c} \bigl(2f_{j-1}(2c+1) - f_{j-1}(2c)\bigr)\bigl(2\chi_{j-1}(2c+1) - \chi_{j-1}(2c)\bigr).
\]
Both tables are then restricted in place (\cref{def:restriction}).
The verifier recovers $s_j(r_j)$ by Lagrange interpolation on $\{0,1,2\}$.

\subsection{Commitments and transcript}\label{sec:impl:merkle}

\paragraph{Committed levels.}
The implementation runs $\Pieval^{\mathcal{J}}$ (\cref{sec:pcs:skip}) with $\mathcal{J} = \{0, k, 2k, \dots\} \cap \iv{0}{\ell-1}$; the parameter $k \in \iv{1}{6}$ is part of the parameters, $k = 4$ by default, and $k = 1$ gives $\Pieval$.
The prover computes every folded word, since each is needed for the next fold, but builds Merkle trees only for $w_j$, $j \in \mathcal{J}$.

\paragraph{Leaves.}
The trees are the salted Merkle trees of \cref{sec:proofs:bcs}: a leaf is $\mathsf{H}_{K_0}(\text{values} \,\|\, \gamma)$ with a uniformly random salt $\gamma$ of $32$ bytes ($= \lambda_{\mathsf{H}}/8$), and an internal node is $\mathsf{H}_{K_1}(\text{left} \,\|\, \text{right})$, where $\mathsf{H}_K$ is BLAKE3 in keyed mode under one of two distinct fixed keys, which separate the two domains; an internal node is then $64$ bytes, one BLAKE3 compression.
The commitment tree has one leaf per fibre of $w_0$, holding its two values, and its leaves are ordered as in \cref{sec:sound:skip}: with the natural ordering of $L$, the fibre over the point of index $f \in \iv{0}{M/2 - 1}$ of $L_1$ is the leaf $t\,2^{d_0-1} + u$, where $f = t + u M/2^{d_0}$ with $t < M/2^{d_0}$, so the $2^{d_0-1}$ fibres of a coset are consecutive leaves.
For $j \in \mathcal{J} \setminus \{0\}$, the tree of $w_j$ has one leaf per coset, holding its $2^{d_j}$ values: leaf $t \in \iv{0}{M_j/2^{d_j} - 1}$ holds $w_j$ at the positions $t + u M_j / 2^{d_j}$, $u \in \iv{0}{2^{d_j} - 1}$.
The salts of a tree are the output of BLAKE3 in keyed mode (as an extendable-output function) under a seed of $32$ bytes drawn from the operating system: one seed per commitment and one per opening, so that the salts are fresh and secret.

\paragraph{Openings and local folds.}
A query point of index $i$ reads, at level $j \in \mathcal{J}$, the coset of index $t = i \bmod (M_j/2^{d_j})$.
For each committed word, the proof opens each distinct leaf once, in increasing order, with its values and salt, and gives the authentication paths of all opened leaves as one merged path: the nodes that cannot be recomputed from the opened leaves, level by level from the leaves to the root, each level from left to right.
The verifier recomputes the set of opened leaves from the query points, and rejects unless the merged path uses every node exactly once and ends at the root.
It then folds each opened coset $d_j$ times (\cref{lem:local-fold}), with the inverses $(2\xi)^{-1}$ obtained from one power of $\omega^{-1}$ per coset and fold step, and compares the result with the entry of the next committed coset, or with $G$.

\paragraph{Fiat--Shamir.}
The transcript is a BLAKE3 hash chain, and follows the rounds of $\Pieval^{\mathcal{J},\Sigma}$ (\cref{sec:sound:pq,sec:sound:skip}).
It absorbs the parameters $(n, R, \ell, \kappa)$, the basis, the salt length, $e$ and $k$, the root of $w_0$, the point $z$ and the value $v$.
In round $j \in \iv{1}{\ell}$ it absorbs $s_j$, then (for $j \ge 2$ and $j - 1 \in \mathcal{J}$) the root of $w_{j-1}$, then a fresh round salt of $32$ bytes, and derives $r_j$.
In round $\ell+1$ it absorbs the final table $g$ and a round salt, and derives the $\kappa$ query points.
A challenge in $\F$ is the image under the map $\phi$ of \cref{sec:sound:pq} of $\beta = 64(e+1)$ output bits ($192$ for $e = 2$, $320$ for $e = 4$): the bits are read as an integer, whose first $e$ digits in base $p$ are the coordinates of $r_j$.
The $\kappa$ query points are $\kappa$ consecutive blocks of $n + R$ output bits, read as exponents of $\omega$ (the map $\mathrm{pos}$).

\paragraph{Relation to the compiler of~\cite{CDHZ25}.}
The implementation has the structure of $\mathsf{BCS}[\Pieval^{\mathcal{J},\Sigma}]$: salted trees, one round salt per round, challenges derived from the roots and salts of all previous rounds, and the maps $\phi$ and $\mathrm{pos}$; merged paths are an encoding of the separate paths (\cref{rem:skip-params}).
The messages that the verifier reads entirely (the round polynomials and $g$) are absorbed into the transcript in the clear rather than committed by a tree, which is the usual encoding of such messages.
With the quartic extension, $\kappa = 248$ and $\beta = 320$, the implementation uses the parameters of \cref{rem:pq-params,rem:skip-params}.

\paragraph{Proof format.}
A proof consists of the $\ell$ round polynomials (three values each), the roots of the words $w_j$, $j \in \mathcal{J} \setminus \{0\}$, the $\ell+1$ round salts, the final table $g$, and, for each committed word, the opened leaves (values and salts) and the merged path.
Its encoding has no length prefixes, since every length is fixed by the parameters and the query points, and the verifier rejects a proof in which any length differs.
The repository contains test vectors: transcript challenges in both extensions, and complete proofs produced from fixed seeds.

\paragraph{Final message.}
The verifier converts $g$ to monomial form with the transform $C^{\otimes (n-\ell)}$ and evaluates $G$ at each $\xi_\ell$ by Horner's rule.
It evaluates $\tilde g(z_{\ell+1},\dots,z_{n})$ for the closure check by $n-\ell$ successive restrictions (\cref{lem:restriction}).

\subsection{Testing}\label{sec:impl:tests}

The test suite checks each component against a direct definition (field arithmetic against $128$-bit integer arithmetic, the transforms against the expansion of \cref{cor:bool-coeff}, the NTT against Horner evaluation) and the protocol end to end:
\begin{itemize}
  \item \emph{Completeness.} Honest proofs verify for both extensions, $k \in \{1,2,3,4\}$, all $1 \le n \le 9$, all $0 \le \ell \le n$ and rates $\rho \in \{1/2, 1/4, 1/8\}$, and the returned value equals $\tilde f(z)$ computed independently.
  \item \emph{Tampering.} Proofs are rejected after changing the claimed value, the point or the commitment; after every single change of one sumcheck value, one entry of $g$, one root, one round salt, one opened value, one leaf salt or one node of a merged path, and after reordering opened leaves, for $k \in \{1,2,3,4\}$; after changing the length of any of these lists or strings; and when a proof for a different table is checked against the original commitment. The verifier returns an error and never panics on such inputs, which randomised truncations also check.
  \item \emph{Fresh salts.} Two commitments to the same table have different roots, and two openings have different round salts; both verify.
  \item \emph{A consistent cheater.} Consider a prover that commits honestly and claims $v+1$. In every round it adds half of the current discrepancy to each of $s_j(0)$, $s_j(1)$ and $s_j(2)$, so that every sumcheck check passes; it folds honestly, and modifies $g(0)$ so that the closure check passes. This prover is rejected in every tested configuration, and always in the query phase, as predicted by case~2 of \cref{thm:sound}.
\end{itemize}
Independently, each algebraic statement of \cref{sec:kernel,sec:fold} and the two lemmas of \cref{sec:sound:lemmas} were checked numerically over small prime fields ($\mathbb{F}_{17}$ exhaustively where feasible, and $\mathbb{F}_{257}$); these scripts are part of the artefact.
